\subsection{Assignment 19 -- Analisi stagionale e per genere dell’impatto dei featuring}

L’obiettivo dell’Assignment 19 è estendere e raffinare quanto analizzato
nell’Assignment 16. Se in precedenza l’attenzione era rivolta all’impatto annuo
degli stream per artista, in questa analisi il focus viene spostato sul ruolo dei
\textit{featuring} all’interno di specifiche stagioni e generi musicali, al fine
di comprendere in quali contesti essi rappresentino un fattore determinante per
il successo commerciale.

Dal punto di vista del business, questa query risulta particolarmente rilevante
poiché consente di passare da un’indicazione aggregata a una strategia mirata.
Non ci si limita a stabilire se le collaborazioni abbiano funzionato in un dato
anno, ma si analizza se il loro impatto vari in funzione della stagione
(ad esempio Estate per le hit stagionali o Inverno) e del genere musicale, nonché
come tali dinamiche si siano evolute dal 1992 a oggi.

La query riprende la logica di ponderazione degli stream introdotta
nell’Assignment 16, assegnando un peso pari all’80\% all’artista principale e al
20\% agli artisti in featuring, introducendo tuttavia alcuni miglioramenti
tecnici per una lettura più accurata dei dati. In particolare:

\begin{itemize}
    \item l’utilizzo di tuple MDX consente di separare in modo esplicito
    l’impatto degli artisti primari da quello dei featuring sulla misura degli
    stream;
    \item i membri \textit{Unknown} sono stati rimossi dalla dimensione
    stagionale tramite la funzione \textit{EXCEPT}, direttamente nel set delle
    righe, garantendo che l’analisi si basi esclusivamente su periodi
    temporalmente definiti;
    \item tramite la funzione \textit{FILTER} sono stati selezionati unicamente
    i casi in cui l’indice calcolato risulta positivo, escludendo gli anni dal
    1992 al 1995, nei quali i featuring erano assenti. Questa scelta rende il
    report più compatto e focalizzato sulle dinamiche effettivamente rilevanti.
\end{itemize}

L’introduzione del dettaglio stagionale consente di individuare con precisione
il momento in cui il fenomeno dei featuring inizia a emergere in modo
significativo. L’analisi mostra infatti che le collaborazioni diventano
rilevanti a partire dall’inverno del 1996; fino a quel momento, il mercato
risulta prevalentemente solista e privo di collaborazioni con un impatto
misurabile.

Il confronto con i risultati dell’Assignment 16 evidenzia ulteriormente il valore
aggiunto di questa analisi. Mentre l’analisi annua precedente forniva una misura
media dell’impatto dei singoli artisti, l’introduzione delle dimensioni
stagionale e di genere rivela come il successo delle collaborazioni sia fortemente
dipendente dal contesto. In stagioni specifiche, generi come Trap e Rap raggiungono
valori dell’indice di convenienza superiori al 20\%, indicando che in tali ambiti
il featuring non rappresenta un evento marginale, ma costituisce il principale
motore degli ascolti stagionali.

Il picco massimo osservato nel 2025, pari al 33.33\% per il genere
\textit{dark\_poetic}, conferma infine come la collaborazione sia diventata una
leva strategica di primaria importanza, in grado di influenzare in modo
significativo il successo stagionale di un intero genere musicale.
